% Gauss-Legendre Quadrature

Gauss-Legendre quadrature approximates the integral
%
\begin{equation}
    \label{eq:sn_legendre_formula}
    \int_{-1}^1 h(u) \, \mathrm{d}u \approx \sum_{i=1}^n w_i h(u_i)
\end{equation}
%
where the abscissas \(u_i\) are the roots of the Legendre polynomial \(P_n(u)\) and the weights \(w_i\) are given by \cite[25.4.29]{Abr72}
%
\begin{equation}
    \label{eq:sn_legendre_weights}
    w_i = \frac{ 2 }{ (1-u_i^2) [P_n'(u_i)]^2 }
\end{equation}
%
The roots and abscissas are usually not calculated but taken from tabulated values, e.g. \cite{Hol10a}.

Adjustment of the integration boundaries may be achieved by a linear transformation of the integration variable:
%
\begin{equation*}
    \int_a^b h(u) \, \mathrm{d}u =
        \frac{ b-a }{ 2 } \int_{-1}^1 h \left( \frac{ (b-a) u + (b+a) }{ 2 } \right)
        \, \mathrm{d}u \approx
        \sum_{i=1}^n \tilde{w}_i \, h \left( \tilde{u} \right)
\end{equation*}
%
with
%
\begin{equation}
    \label{eq:sn_legendre_transform}
    \tilde{u}_i = \frac{ (b-a) u_i + (b+a) }{ 2 } , \qquad
        \tilde{w}_i = \frac{ b-a }{ 2 } w_i ~.
\end{equation}
%
Gauss-Legendre quadrature integrates polynomials of order \(< 2n\) exactly, so it is most accurate if all derivatives up to order \(2n-1\) exist and higher derivatives are small.
